% Generated draft: requires mathematical and source review.
\documentclass[11pt]{article}
\usepackage[margin=25mm]{geometry}
\usepackage{fontspec}
\setmainfont{Latin Modern Roman}
\usepackage{amsmath,amssymb,mathtools,unicode-math}
\setmathfont{Latin Modern Math}
\usepackage{enumitem,needspace,xurl,hyperref}
\setlength{\parindent}{0pt}
\setlength{\parskip}{5pt}
\setlength{\emergencystretch}{3em}
\newcommand{\N}{\mathbb{N}}
\newcommand{\Z}{\mathbb{Z}}
\newcommand{\Q}{\mathbb{Q}}
\newcommand{\R}{\mathbb{R}}
\newcommand{\C}{\mathbb{C}}
\newcommand{\F}{\mathbb{F}}
\newcommand{\A}{\mathbb{A}}
\newcommand{\PP}{\mathbb{P}}
\newcommand{\E}{\mathbb{E}}
\newcommand{\eps}{\varepsilon}
\newcommand{\dd}{\,\mathrm{d}}
\begin{document}
\section*{UnsolvedMath ID: 11000074}
\textbf{Source problem code:} catalogue entry
\subsection{Definitions and mathematical statement}
Let $\Sigma_g$ be a closed oriented genus-$g$ surface, let $\operatorname{Homeo}^+(\Sigma_g)$ be its orientation-preserving homeomorphism group, and let
\[
 \pi:\operatorname{Homeo}^+(\Sigma_g)\longrightarrow\operatorname{Mod}_g
\]
send a homeomorphism to its isotopy class. For $H\le\operatorname{Mod}_g$, a section over $H$ is a group homomorphism $s:H\to\operatorname{Homeo}^+(\Sigma_g)$ with $\pi\circ s$ equal to the inclusion of $H$ in $\operatorname{Mod}_g$.

\textbf{Problem 6.5 (the sections problem).} Determine exactly which subgroups $H$ admit such a section. Ask the same question with the target replaced by specified regularity or geometric subgroups: $\operatorname{Diff}^r(\Sigma_g)$ for $r=1,2,\ldots,\infty,\omega$, area-preserving diffeomorphisms, quasiconformal homeomorphisms, and similar groups. The diffeomorphism targets are understood in the orientation-preserving component convention appropriate to the projection; $r=\omega$ means real analytic. Area preservation requires a fixed area form, and quasiconformality a fixed conformal structure, to be supplied. No uniform quasiconformal distortion bound is imposed by the excerpt. For each target $A$, the equation is $\pi_A\circ s=\operatorname{incl}_H$ with $s:H\to A$.

The source records the following context, with theorem hypotheses and present status requiring verification:
\begin{enumerate}
\item Sections exist for free $H$ in the homeomorphism setting, by choosing representatives of free generators.
\item For free abelian $H$, sections exist even into $\operatorname{Diff}^{\infty}(\Sigma_g)$, using the abelian-subgroup classification of Birman--Lubotzky--McCarthy [BLM].
\item Finite $H$ admit real-analytic sections by Kerckhoff's Nielsen realization theorem [Ke], realized as automorphisms of a suitable Riemann surface.
\item Morita [Mo5] is reported to obstruct a $C^2$ section over the whole group when $g\ge5$, via a codimension-two foliation of the universal curve, Bott vanishing, and a nonzero Miller--Morita--Mumford class. The extraction writes the target as $\operatorname{Diff}^2(\Sigma_2)$ in this assertion, which conflicts with the genus-$g$ projection; this apparent subscript error must be checked, not silently accepted as a theorem about genus two. The suggested extension of this argument to finite-index subgroups is described as prospective.
\item Markovic [Mar] is reported to obstruct a section over $H=\operatorname{Mod}_g$ even into homeomorphisms. The excerpt does not give precise genus hypotheses. It stresses that torsion and braid relations used in that argument do not supply the desired result for general finite-index subgroups.
\end{enumerate}
The classification therefore includes the finite-index-subgroup issue highlighted by the source. Its final reference to a following question does not supply a further separate statement here. The map $\pi$ in this problem is the homeomorphism-to-mapping-class projection, not the forgetful map in the Birman exact sequence.
\subsection{Short English statement}
Classify mapping-class subgroups that can be realized homomorphically by homeomorphisms, and by representatives of specified smoothness or geometric type. Retain the finite-index-subgroup problem and check the genus typo in the quoted $C^2$ obstruction.
\subsection{Sources}
\begin{itemize}
\item Local UnsolvedMath registry, permanent ID 11000074, source code catalogue entry; no individual proposer or proposal year supplied.
\item Farb, edited volume \emph{Problems on Mapping Class Groups and Related Topics}, Problem 6.5, PDF page 55; \url{https://www.math.uchicago.edu/~farb/papers/mcgbook.pdf}. The registry records PDF-font-structure extraction and access on 2026-07-29; the inventory distinguishes 329 labeled headings by source page.
\item Source labels [BLM] (Birman--Lubotzky--McCarthy), [Ke] (Kerckhoff), [Mo5] (Morita), and [Mar] (Markovic). Their complete bibliography entries and theorem hypotheses are absent from this record.
\end{itemize}
\end{document}
